Digital services are often perceived as immaterial. The "Cloud" suggests something ethereal and weightless. This perception obscures the physical reality. As Crawford argues, artificial intelligence and digital infrastructure are "profoundly material," requiring minerals extracted from the earth, electricity generated by power plants, and water consumed for cooling \cite{Crawford2021Atlas}. When designers make interface decisions, they are making claims on these material resources. This chapter highlights some UX decisions. Firstly the section highlights the power of defaults with the example of video streaming. Then it showcases three UX pattern examples that have eco-friendly alternatives. Lastly it talks about AI as a trend and its problems.

\subsection{Defaults and Streaming Quality}

When designing a video stream media player a designer must choose a default resolution. Default settings and digital nudging, which is the use of user-interface design elements
to guide people’s behavior in digital choice environments \cite{nudge},
 are very powerful. In 2011 Jared Spool analyzed hundreds of Microsoft Word users settings files and found that less than 5\% changed any of the 150+ settings at all \cite{Spool2011Settings}. Similarly in 2003 Eric J. Johnson and Daniel Goldstein showed that defaults have a big impact for even something very personal such as organ donation. In a country where opt-in is needed it resulted in 12\% organ donors, whereas a country with opt-out, like Austria, being an organ donor had an acceptance rate of 99.98\% \cite{Organs}. Also with a direct connection to environmental footprint and responsibility, research showed that when a cargo transportation platform changed the default option from opt-in to opt-out for CO₂ offsetting, participation changed from 12\% to stabilized 50\% after 12 weeks. Additionally the firm didn’t lose revenue \cite{Transports}.

\par Video quality defaults are a possible UX decision that has an environmental impact. This choice determines energy consumption for users, who will likely never adjust the setting or opt-out from the default. In 2018, researchers measured the same device on the same network and found that streaming in HD produced 130 grams of CO₂ per hour, while 4K produced 610 grams of CO₂ per hour which is nearly five times more. In 2024 a study by the Carbon Trust that found an entire European average of 55 grams per hour. This shows a general problem for these statistics. They are very context heavy and are dependent on energy grid, network and video quality, but the directional effect is consistent: higher resolution results in higher energy consumption \cite{2018Report, CarbonTrust2021Streaming}. For a single user this seems insignificant but for Netflix's 325 million subscribers streaming 95 billion hours in the first half of 2025 alone \cite{Netflix2025H1, Netflix2025Q4}, even small per-user differences aggregate to large numbers. This shows individual design choices, multiplied across hundreds of millions of users who rely on default settings can result in significant environmental impact.

\subsection{UX Pattern Decisions}

Designers exercise environmental influence through everyday pattern decisions that dictate the "material weight" of digital interfaces. This is showcased by following examples:
\begin{itemize}
  \item \textbf{Autoplay:}
  While implemented to maintain engagement, empirical research demonstrates that autoplayed preview videos can increase data transfer by up to 86\% compared to static interfaces. This creates a large carbon footprint increase as platforms decode high-bit-rate video. \cite{Andreasson2024Autoplay}.
\item \textbf{Infinite Scroll:}
Infinite scroll is classified as a "dark pattern" that exploits psychological vulnerabilities, similar to slot-machines, to bypass user intent and maximize attention capture\cite{Roffarello2022DarkPatterns}. By removing natural stopping points, this design forces the continuous loading of data centers and network resources for content that may never be actively consumed.
\item \textbf{Media Assets:}
A paper in 2023 analyzed two websites that were hosted with the same method. The difference was one site was a simple text-only website with little interaction. It generated 0.03 g of CO₂ per visit. The other site was more complex, designed with images, animated gifs, embedded videos and written articles. The second website generated 0.92 g of CO₂ per visit, which shows the big impact media assets can have \cite{Vacanti2024}.
\end{itemize}
These decisions again suffer from the scaling effect. One page visit seems negligible, but when a website gets millions of daily views or millions of social media users scroll longer than they actively wanted the environmental impact scales up.


\subsection{AI as a Trend}

AI reveals both the scale of impact and the opacity problem. It's estimated that training GPT-4 consumed 50 gigawatt-hours of energy. That is enough to power a city like San Francisco for three days. Training is a one-time cost. What matters more is inference, every time someone queries these models. Inference now accounts for 80 to 90\% of AI's total computing power \cite{MIT}.

\par In August 2025, Google published the first comprehensive, first-party measurement of production AI serving energy consumption. Using a methodology that accounts for active AI accelerators, CPU and DRAM, idle machines, and data center overhead, they measured the median Gemini Apps text prompt at 0.24 watt-hours, emitting 0.03 grams of CO₂ equivalent and consuming about 0.26 milliliters of water. This is roughly the energy impact of watching TV for about 9 seconds. Google claims this value has gotten about 44x more efficient over a 12 month period \cite{google}.

\par However, before Google's paper researchers argued that closed AI model providers are serving up a total black box . This is still mostly true and Google also didn't disclose data on complex prompts, image generations and video generation. This is a lack of transparency that is increasingly dangerous as we transition toward "reasoning models" which require up to 43 times more energy for simple tasks than traditional models. Additionally inference aggregates into a surge in infrastructure demand. After a decade of relative stability, data center electricity consumption has doubled since 2017. Projections suggest that by 2028, data centers will claim 12\% of total U.S. electricity. Because AI requires "24-7" constant power, these facilities are increasingly turning to dirtier, carbon-intensive energy sources to meet immediate needs, resulting in a carbon intensity 48\% higher than the U.S. average \cite{MIT}. This showcases that AI as a trend has a significant environmental footprint and that the whole "Black box" dynamic brings further problems, when trying to make informed decisions.

\par These examples demonstrate that UX design decisions have measurable environmental
consequences that scale across billions of users. This raises an ethical question:
who bears responsibility for these consequences, and what principles should guide
design decisions? Section 3 addresses this through frameworks of climate justice
and autonomy.
